\section{Methods}

\subsection{Workflow of model evaluation}
The study has three parts that correspond to the three questions. (Q1) Published numbers are audited (the paper audit and the reproducibility funnel), released checkpoints are scored on official and independent test sets, and the authors' pipelines are rerun (author-recipe reruns). (Q2) Models with complete training code are retrained under one protocol on IDRiD and on DDR (unified retraining), including nnU-Net through its own entry point and foundation models under four kinds of use; all checkpoints are applied to six independent test sets (independent tests). (Q3) One factor is varied at a time on representative models (U-Net, nnU-Net, M2MRF, SAM2-UNet, RETFound with a UPerNet head): training source, annotation granularity, input resolution, training amount and loss; the HRDecoder case study of the main text applies the same machinery to one model's own geometry. Computational cost is recorded for every run. Each experiment is named in the text by what it does; the identifiers under which the runs, the results table and the planning documents of the repository file them are listed in Supplementary Table~18. The evaluation protocol and all splits were committed to the repository before the first main-table run. Pipeline checks had scored a reference U-Net and the released M2MRF checkpoints on the IDRiD and DDR test sets before that; the field-of-view rule and the lesion-size boundaries were then revised once, from training images only. After the freeze, the training amount and the selection rule were set from training loss and validation curves, and no test metric was used. One exception is disclosed in Supplementary Note~2: the source recipe of the reference U-Net was revised after the test score of a smoke run had been seen; this affects the source-recipe track only.

\subsection{Data collection}
\paragraph{Training sources.} IDRiD\cite{porwal2018idrid} provides 81 images of $4288\times2848$ pixels with pixel masks for MA, HE, EX and SE; the official split is 54 training and 27 test images. A missing mask file means that the image has no lesion of that class. DDR\cite{li2019ddr} provides 757 images with the official split 383\,/\,149\,/\,225.

\paragraph{Independent test sources.} TJDR\cite{mao2023tjdr} contains 561 images from one hospital: 257 standard-field images ($2048\times2048$, 35--50$^\circ$) and 304 ultra-widefield images ($3912\times3912$, 133$^\circ$), annotated with polygons by two ophthalmologists independently and adjudicated by a third; the four classes are mutually exclusive. The two fields are treated as two datasets. MAPLES-DR\cite{maplesdr2024scientificdata} adds pixel labels for 198 MESSIDOR\cite{decenciere2014messidor} images (45$^\circ$; $1440\times960$ to $2304\times1536$ pixels): seven retinologists drew MA, HE, EX and SE at $1500\times1500$ on the field-of-view crop, MA, HE and EX starting from a network's pre-segmentation and SE from scratch; the labels are placed back on the original image, and only the official test split of 60 images is used. The FGADR Seg-set\cite{zhou2021fgadr} contains 1,842 images of $1280\times1280$ pixels with masks for the four classes and, as extension labels, intraretinal microvascular abnormalities (159 images) and neovascularization (49). Retinal-Lesions\cite{wei2020retinallesionslesionnet} (release 20191227) contains 1,593 images from EyePACS, cropped to the field of view and resized to $896\times896$, from 1,505 patients, annotated with ellipses and polygons by at least three of 45 graders. DiaRetDB1 v2.1\cite{kauppi2007diaretdb1} contains 89 images of $1500\times1152$ pixels (50$^\circ$) in which four experts independently marked regions (4,135 circles, 763 ellipses, 1,036 polygons), each with a confidence level.

\paragraph{Granularity tiers.} Sources are assigned to four tiers: fine (IDRiD, DDR, TJDR standard field, TJDR ultra-widefield, MAPLES-DR), mixed (FGADR), coarse (Retinal-Lesions) and region consensus (DiaRetDB1). Results are averaged only within a tier.

\paragraph{Label mapping.} Only the four classes are mapped; other classes are stored as extension labels and do not enter the four-class mean. TJDR: palette 1, 2, 3, 4 to EX, HE, MA, SE. FGADR: HardExudate, Hemohedge, Microaneurysms, SoftExudate. Retinal-Lesions: hard\_exudate, retinal\_hemorrhage, microaneurysm, cotton\_wool\_spots; vitreous and preretinal hemorrhage are not merged into HE. DiaRetDB1: Hard\_exudates, Haemorrhages, Red\_small\_dots, Soft\_exudates. MAPLES-DR: Microaneurysms, Hemorrhages, Exudates, CottonWoolSpots; its uncertain red lesions (MA or HE) and uncertain bright lesions (EX, SE or drusen) are ignored for the two classes they could belong to, and drusen and neovascularization are extension labels.

\subsection{Model collection}
The literature search was frozen at 30 September 2026. A fixed query (three concepts in title or abstract: diabetic retinopathy; lesion, microaneurysm, exudate, hemorrhage or cotton wool; segmentation) was run on PubMed, Europe PMC and OpenAlex for 2015 onwards. It identified 1,389 records, of which 948 passed the Boolean gate and 764 remained after de-duplication by title; 112 entered the full-text queue. Screening of a priority list (27 DOIs, one dataset lead and one further paper found when a mis-assigned DOI was corrected) required colour fundus photographs, pixel-level segmentation of at least two of the four classes and an obtainable public dataset; it included 24 methods and 2 datasets and excluded 3 records. Together with a snowballed corpus the ledger holds 86 records. Only the authors' own repositories count as official code.

\rn{n.models} model variants were registered in three tiers (\rn{n.tier1}, \rn{n.tier2} and \rn{n.tier3}). Each has a reproducibility level: R0, no code; R1, model definition only; R2, inference code and weights; R3, training code; R4, training code, weights, splits and environment. Each is assigned to evaluation tracks: A, author weights; B, unified retraining of the official code; C, re-implementation from the paper, marked in the tables; a re-implementation fixes every unstated choice in a written list of assumptions before any experiment and is not tuned to the published numbers, and its parameter count is compared with the published one and reported. Methods whose repository holds the model but no training loop (WSRFNet) or whose code cannot run on current hardware (PMCNet, TensorFlow 1.13) are retrained under Track B with the loss and optimizer the paper states, written into the wrapper and disclosed as such. One method with official code and weights, HRDecoder, was missed by the frozen search and added on 14 October 2026 after a targeted check for code-available methods; the addition is recorded as a protocol deviation. Authors of methods without code received a fixed e-mail on 3 October 2026 asking for code, weights and splits, with a deadline of four weeks; where code exists without weights, the weights are requested separately. Non-responses and failures remain in the denominator of the reproducibility funnel.

\subsection{Preprocessing}
\paragraph{Unified layout.} Every source is exported to the same layout: the image at original resolution with the bytes of the release, one binary mask per class, a field-of-view mask and the split lists. A missing mask file is an all-negative image for that class; the image stays in the split.

\paragraph{Field of view.} The field-of-view (FOV) mask is the largest connected component of pixels with luminance $0.299R+0.587G+0.114B>10$, with holes filled. An Otsu threshold was abandoned after it removed the darker half of the retina in several training images. With the fixed threshold, the largest share of lesion pixels lying more than 2 pixels outside the mask on the training images is $2.3\times10^{-5}$ (IDRiD hemorrhages) and 0 on DDR. For the independent sources one rule is added: if the four $40\times40$ corner blocks are uniform (1st to 99th percentile within 5 grey levels), the threshold is the larger of 10 and the corner median plus 3. It changes one of 4,085 images, a DiaRetDB1 image with a uniform grey background.

\paragraph{Canvas.} For training and inference the image is cropped to the FOV and resized so that its long side is 1,440 pixels. Predictions are mapped back to the original resolution by bilinear interpolation, set to zero outside the FOV and stored as 16-bit images (probability = value\,/\,65535). All metrics are computed offline from these maps.

\paragraph{Source-specific handling.} FGADR masks have anti-aliased edges with values 1--254; a pixel is positive at $\geq128$, and the number of pixels dropped is recorded per class. In Retinal-Lesions, value 255 is lesion and value 127, which the authors use for annotations that do not follow the grading guideline, is written to a per-class ignore region that is skipped by all metrics and, when this source is used for training, by the loss. In DiaRetDB1, circles, ellipses and polygons are rasterized at pixel centres; for one expert and class the highest confidence per pixel is kept, with High = 1.0, Medium = 0.75 and Low = 0.25; the four experts are averaged and a pixel is positive at $\geq0.75$, the threshold of the original protocol. With this rule six of the eight positive-image counts of the original paper are matched exactly and two differ by one image (the preparation report in the repository); the union of all expert marks is kept as a lenient ground truth for sensitivity analysis.

\paragraph{Exclusive training targets.} Models with a softmax head need one label per pixel. Overlapping pixels are written in the order EX, HE, SE, MA, later classes overwriting earlier ones. This removes less than 0.5\% of the positive pixels of any class in the training sets (at most 0.34\% at original resolution). Evaluation always uses the four original, possibly overlapping masks.

\subsection{Construction of training, validation and test sets}
Official splits are never changed. New splits use the generator PCG64 with seed 20261003, and the lists were committed before any training on the source concerned. Where a split is stratified, quotas per stratum follow the largest-remainder method. IDRiD: 10 of the 54 official training images, drawn at random from the sorted identifiers, are held out for validation (44\,/\,10\,/\,27); author-recipe retraining uses all 54. DDR: official 383\,/\,149\,/\,225. TJDR: the official test set is kept and 15\% of the official training images of each field are drawn for validation, giving 172\,/\,30\,/\,55 (standard) and 209\,/\,37\,/\,58 (ultra-widefield). FGADR: stratified by file-name suffix and DR grade, 30\% test and 15\% of the remainder for validation (1,096\,/\,193\,/\,553). Retinal-Lesions: split by patient, stratified by the higher grade of the two eyes (951\,/\,164\,/\,478). DiaRetDB1: official 28\,/\,61; the 28 training images are used only to choose thresholds.

\subsection{Test scenarios and leakage audit}
The internal test is the official test set of the training source. Independent tests come from sources excluded as a whole from training and from internal testing; conclusions rest on them, as in the template study, where independent databases were held out and similar sequences were removed\cite{lu2026tcrepitopenaturemethods}. Every main-table checkpoint is applied to the test splits of the five independent sources. Ultra-widefield images are predicted at their native resolution with sliding windows, because matching pixels per degree to the 1,440-pixel canvas would require upsampling; a version resized to 1,440 pixels is reported for comparison. The per-degree scale is defined by the field of view of the training source; for a run whose training sources do not fix one (the factor cells trained on the standard-field TJDR images, on FGADR, or on merged sources), only the 1,440-pixel version is run and stands for the ultra-widefield result. For author weights, a source that the original model was trained on is removed from its independent tests. Near-duplicate images are searched with perceptual hashes and feature similarity between sources, between the splits of a source, and between the pretraining sets of foundation models and the test sets; an independent test image that duplicates a training image is removed and listed (Extended Data Fig.~3).

\subsection{Evaluation of the original models}
\paragraph{Released checkpoints (author weights).} Released checkpoints are run in a current framework; for M2MRF, whose repository requires a framework version that does not support the cards used here, the forward pass was checked by reproducing the repository numbers with the authors' function. Each checkpoint is reported with three numbers: the paper, the authors' evaluation function, and this study's protocol. The authors' function for M2MRF evaluates the whole image, overwrites overlapping labels in the order EX, HE, SE, MA and computes AUPR from 11 thresholds.

\paragraph{Author-recipe reruns.} M2MRF and HRDecoder are retrained with the authors' pipelines, augmentation and training set (all 54 IDRiD training images; three M2MRF seeds and one HRDecoder run, plus a first M2MRF run with the wrong effective batch that is reported separately) and scored at the last iteration without checkpoint selection, with the authors' function and with this study's protocol. The registered plan had a DDR run and HACDR-Net's recipe as well; these were cut when the scope was fixed, and the second and third M2MRF seeds were run when cards were free. For M2MRF: binary Dice loss (smoothing $10^{-5}$); SGD with learning rate 0.01, momentum 0.9 and weight decay $5\times10^{-4}$; polynomial decay with power 0.9; 40,000 iterations on IDRiD and 60,000 on DDR; random crops of $960\times1440$ after resizing; photometric distortion; HRNet-W48 initialized from ImageNet. For HRDecoder: the authors' segmentor and loss, images resized to $2880\times1920$ with scale jitter 0.5--2, random $1920\times1920$ crops, flips and rotations of $\pm45^\circ$, SGD 0.01 with 1,500 warm-up steps and linear decay over 20,000 steps, four images per step; their four-process synchronized batch normalization becomes one process with four images. The authors' runner evaluates on the test split every 2,000 iterations and keeps the best checkpoint; the rerun reports the last step, and what that test-split selection would have returned is given separately.

\subsection{Unified retraining}
\paragraph{What is shared.} Data, splits, the 1,440-pixel canvas, augmentation, the training amount in pixels, checkpoint selection, inference geometry and evaluation are the same for every model. Each model keeps its loss, optimizer, learning rate, weight decay, schedule shape and effective batch from its source recipe (Table~\ref{tab:recipes}); the schedule is stretched over the shared number of steps.

\paragraph{Training amount.} One equivalent epoch draws as many patch pixels as the canvases of the training split contain, so that a model trained on $224\times224$ patches draws about 18 times more patches per epoch than one trained on $960\times960$ patches:
\begin{equation}
\text{steps per epoch}=\left\lceil \frac{1}{B}\left\lceil \frac{\sum_i H_i W_i}{h\,w} \right\rceil \right\rceil ,
\end{equation}
where $H_i\times W_i$ is the canvas of training image $i$, $h\times w$ the patch and $B$ the effective batch. The number of equivalent epochs was fixed per dataset by a budget probe: U-Net and M2MRF were trained once (seed 0) at twice a provisional amount, and only training loss and validation \mAUPR{} were read. The validation curve was smoothed with a centred three-point mean; if the smoothed maximum of any probe of a dataset lay in the last 10\% of the steps, the amount was doubled. This gave 1,000 equivalent epochs for IDRiD and 400 for DDR. The rule was revised twice after the validation curves had been seen and never after a test metric had been read; all versions of the rule and their outcomes are given in Supplementary Note~2.

\paragraph{Patches and inference.} Models with a fixed input size (Swin-Unet and RETFound with UPerNet, 224; FCT, 384; CLC-Net, TC-Net and SAMed, 512; SAM-Adapter and MedSAM, 1,024) are trained on patches of that size cut from the canvas and predict with sliding windows of the same size at an overlap of 0.5; HRDecoder keeps its authors' windowed inference at 960. All other models are trained on $960\times960$ patches and predict on the whole canvas (Table~\ref{tab:recipes}).

\paragraph{Augmentation.} Random scaling by 0.5--2 before cropping, horizontal flip with probability 0.5, rotation by multiples of 90$^\circ$, and photometric distortion with probability 0.5 (brightness $\pm32$, contrast and saturation 0.5--1.5, hue $\pm18^\circ$). The randomness of a sample depends only on the seed and the sample index.

\paragraph{Checkpoint selection.} A run is validated 20 times at equal intervals and the checkpoint with the highest validation \mAUPR{} inside the FOV is scored; the last step is kept for a sensitivity analysis. A model whose best checkpoint falls in the last two validations is flagged as possibly under-trained. The test split is not read during training.

\paragraph{Reproducibility.} A resume file stores model, optimizer, step, selection state, random-number state and data order; a run resumed after an interruption is bit-identical to an uninterrupted one. Each run writes its environment (host, card, driver, framework, seed, commit at process start, parameter count, wall time) and its recipe with every deviation.

\paragraph{Rare-class collapse.} The reference U-Net uses a binary cross-entropy plus a sigmoid Dice term computed per image and class with smoothing 1. On a patch without positives this term is minimized by an all-zero output with a gradient of about 1 per pixel, whereas on a patch with a few positive pixels the gradient per pixel is about 2\,/\,(number of positive pixels + 1). Under patch training 83\% of DDR patches and 67\% of IDRiD patches contain no soft exudate. In the budget probe the soft exudate channel collapsed to zero on IDRiD, and on DDR the microaneurysm, hemorrhage and soft exudate channels did. With a smoothing of $10^{-5}$ the gradient of this term on an empty patch is small while the channel's output is far from zero, and a Dice pooled over the batch (nnU-Net) is not affected. The U-Net's Dice term was therefore pooled over the batch per class with smoothing $10^{-5}$, as in nnU-Net; the revision was decided from the training and validation curves of a diagnostic run (validation \mAUPR{} 0.569 on IDRiD and 0.544 on DDR against 0.127 before) after the registered recipe's test numbers had been read, and both versions are reported, the registered one in Supplementary Table~6 (Supplementary Note~2). Before any foundation model was trained the same rule was registered for SAM2-UNet, whose weighted IoU term is computed per image and class with smoothing 1, and for MedSAM, whose Dice term is computed per image and class with squared predictions and smoothing $10^{-5}$: if the validation AUPR of a rare class stays at its prior level throughout training, the term is pooled over the batch and the model is retrained from its initial weights. The rule applied to SAM2-UNet on DDR (microaneurysm and soft-exudate validation AUPR 0.02--0.04 and 0.05--0.10 from the first to the last evaluation) and to MedSAM on DDR (microaneurysm 0.000--0.004 and soft exudate 0.002 at every evaluation, their prior levels, while hemorrhage and exudate trained); both models were retrained with the term pooled over the batch (SAM2-UNet to 0.664 and 0.360 \mAUPR{} on IDRiD and DDR, MedSAM to 0.642 and 0.413), the main tables show the retrained runs, and the collapsed runs are kept as the registered recipe in the supplementary tables.

\paragraph{nnU-Net.} nnU-Net is run through its own entry point in 2D at original resolution, with the same splits written into its split file and exclusive labels composed in the same overwrite order; softmax probabilities are exported, mapped to the original image and scored with the same evaluation. The default configuration and the residual-encoder preset M are both run. Its preprocessing is its own and is marked as such.

\paragraph{Source-recipe track.} A second track, each model trained with its source recipe after minimal adaptation to the shared data, canvas and evaluation (seed 0), was registered to measure what moving a recipe to this task gives; it is not an architecture ranking, since the recipes differ about 67-fold in the number of images seen on IDRiD. Six of its runs were started and the track was stopped for compute on 2026-10-11 before any of them finished; no number from it is reported.

\subsection{Foundation models}
Four kinds of use are reported separately. (a) Zero-shot: automatic masks of SAM and SAM~2, scored once as class-agnostic ``any lesion'' proposals and once with each mask assigned the class it overlaps most; text prompts for SAM~3 with three wordings per class. (b) Oracle prompts: boxes or points generated from ground-truth connected components, with 1, 5, 10, 20 or all components per image and class; MedSAM, which accepts boxes only, belongs here. (c) Parameter-efficient fine-tuning: SAMed, SAM-Adapter and SAM2-UNet, trained on patches cut from the canvas (512, 1,024 and 960 pixels; Table~\ref{tab:recipes}). (d) Encoder plus decoder: MedSAM fully fine-tuned, and RETFound with a UPerNet head. To isolate the effect of ophthalmic pretraining, two pairs of ViT-L encoders fixed before any was trained are run on DDR with the same UPerNet head, recipe and training amount: RETFound's masked-autoencoder encoder against the ImageNet masked autoencoder it was initialised from, and RETFound's DINOv2 encoder against DINOv2 itself (one seed each; the pair differences are the pre-specified comparisons). After the DDR result of the first pair had been read, the three control encoders were also trained on IDRiD under the same head and recipe, so that both pairs exist on both training sources; this addition is reported as such. Settings (c) and (d) follow the unified protocol.

\subsection{Evaluation of factors}
Each experiment changes one factor and keeps the unified protocol otherwise; every cell is trained at the pixel budget of DDR with the DDR recipe of its model and one seed (seed 0), and differences smaller than twice the median seed standard deviation of the main table are reported without interpretation. \emph{Training source}: the five representative models (U-Net, nnU-Net, M2MRF, SAM2-UNet and RETFound with UPerNet) are trained on the standard-field TJDR images and on FGADR and tested on all eight test sets, which with the main table gives four training sources. \emph{Annotation granularity}: U-Net and M2MRF on six training sets at an equal number of patch pixels, namely DDR, FGADR, Retinal-Lesions, DDR with FGADR, DDR with Retinal-Lesions, and DDR with 383 FGADR images drawn once (list committed before training); the DDR and FGADR rows are the main-table and training-source runs. \emph{Input resolution}: M2MRF and RETFound with UPerNet on DDR with a long side of 512 or 1,440 pixels (the main table) and with patches at native resolution; the HRDecoder case study of the main text trains that model at its authors' 2,880-pixel canvas and 1,920-pixel window with the equivalent epochs of the unified protocol counted at that canvas. \emph{Training amount}: U-Net and M2MRF on 10\%, 50\% and 100\% of the DDR training set (nested subsets drawn once) and on IDRiD, DDR and standard-field TJDR merged. \emph{Loss}: U-Net on DDR with cross-entropy, batch Dice, their sum (the main table) and the per-image Dice of the registered recipe, with the architecture fixed. \emph{Learning rate}: HRNet-W48 and M2MRF on DDR with the learning rate multiplied by 1/3 and 3, judged on validation \mAUPR{} only. An extra vessel-map input and further resolution and training-amount levels of the registered plan were cut for compute and are listed as future work.

\subsection{Metrics}
\paragraph{Pixel level.} For each class, the positive and negative pixels of all test images inside the FOV (and outside the ignore region of that class) are pooled, and AUPR is computed exactly over the 16-bit probability levels of the stored maps; this equals the average precision
\begin{equation}
\mathrm{AUPR}=\sum_{k}\left(R_k-R_{k-1}\right)P_k ,
\end{equation}
with precision $P_k$ and recall $R_k$ at the $k$-th distinct threshold. \mAUPR{} is the mean over the four classes. Dice and intersection over union are reported at a fixed threshold of 0.5 and at the per-class threshold that maximizes Dice on the validation set. For an image without positive pixels of a class the per-image Dice is missing, not zero.

\paragraph{Lesion level.} Lesions are 8-connected components. A ground-truth lesion is detected if any predicted positive pixel overlaps it; a predicted component is a true positive if it overlaps any ground-truth pixel of the class. Precision, recall and F1 are pooled over the split. Lesion size is the equivalent diameter of the ground-truth component divided by the equivalent diameter of the FOV; the boundaries between small, medium and large, 0.005 and 0.011, are the medians of microaneurysms and hemorrhages on the 44 IDRiD and 383 DDR training images, rounded up to 0.001. Validation and test images were not used to set them.

\paragraph{Annotation style.} The labelling style of each source (Fig.~\ref{fig:sources}a) is measured on the prepared masks of all its images, every image once whatever its split: per class, the 8-connected components are counted and their equivalent diameter is divided by the long side of the FOV bounding box (which for a circular field equals the FOV diameter); ignore regions are removed before counting, and DiaRetDB1 is measured on the $\geq0.75$ consensus the protocol evaluates. Reported are the mean number of components per image, over all images of the source, and the median relative diameter over all components.

\paragraph{DiaRetDB1.} The primary result follows the official image-level protocol: the image score of a class is the maximum probability inside the FOV, and the area under the ROC curve and the weighted error rates are reported. Pixel-level AUPR against the consensus regions is given separately and marked as region ground truth.

\paragraph{Authors' implementation.} The 11-threshold AUPR is reported next to the exact value for all models, only to explain distances to published numbers.

\paragraph{Case study of one model's geometry.} HRDecoder trains its backbone at half the resolution of a crop that is itself two thirds of the authors' 2,880-pixel canvas; the unified protocol halves the canvas and keeps that ratio. Three settings separate the recipe from the budget and the canvas: the authors' recipe retrained (the author-recipe rerun), the unified protocol at twice its training amount (for HRDecoder and, as a control, HRNet-W48 on DDR, trained on cards taken out of the queue), and the unified protocol on the authors' canvas with the authors' 1,920-pixel window and equivalent epochs counted at that canvas (four times the pixels per step). All three were registered before any of them was trained (docs/paper/14, note 46 of the repository).

\paragraph{Ensembles.} Probability maps of the main-table models trained on one source are averaged pixel by pixel in five compositions fixed before any ensemble was scored: all models; the three and the five best by validation \mAUPR{}; the best model of each family; and each family on its own. The member list is the 18 main-table rows scored when the compositions were registered; MedSAM, scored afterwards, is not a member. Members are chosen on validation data only; the ensemble's own validation maps set its thresholds; the test sets and the scoring are those of the single models. Ensembles are reported as a reference below the single models and do not enter the rank tests.

\paragraph{Reference model.} As a starting point for users of the benchmark, two models are trained once on the merged training splits of the three fine-grained sources (IDRiD, DDR and the standard-field TJDR set; 599 images) with the recipe, canvas and pixel budget of their DDR main-table run, validated on DDR, and their probability maps are averaged as in the ensembles. The two were chosen before any merged run by their DDR validation \mAUPR{} among the main-table models (RETFound with UPerNet, 0.643; SAMed, 0.632); the pair, the budget rule and the validation source were registered on 2026-10-23. The reference model is scored on every test set like the ensembles (thresholds from the DDR validation split of the averaged maps), is not a ranked row, and is released with its weights and maps and with a command that applies it to new images (\texttt{python -m bench.reference predict}: field-of-view detection by the rule above, the 1,440-pixel canvas, each member's sliding-window inference, the averaged probability map at the original resolution, and masks at the validation thresholds).

\paragraph{Human reference.} On DiaRetDB1 each of the four experts in turn is scored as if it were a model: the expert's own markings, with the confidence mapping above, are the probability map, and the consensus of the other three (mean confidence $\geq0.75$) is the ground truth; thresholds come from the official training list as for the models, and the four values of each metric are reported as a band next to the DiaRetDB1 column. A map with four confidence levels gives a four-point precision--recall curve, so this band is a lower bound on what an expert's pixel-level AUPR would be with a continuous score. Lesion-level markings by three ophthalmologists on DDR and MAPLES-DR test images are pre-registered and will be reported in a Supplementary Note when complete.

\paragraph{Results by DR grade.} For the test sets that carry a DR grade per image (DDR, FGADR, Retinal-Lesions, MAPLES-DR), every score is also computed within grades 0--1 and within grades 2--4 from the stored per-image histograms, with the thresholds of the full test; ungradable images enter neither stratum. Image-level detection summarizes the maximum probability inside the FOV against the presence of the class in the ground truth.

\paragraph{Label revision.} Refined IDRiD\cite{chankhachon2026refinedidrid} re-annotates the IDRiD images as one $1024\times1024$ grey-level label per image (one class per pixel: background, vitreous hemorrhage, retina, fovea, vessels, optic disc, hard exudates, IRMA, hemorrhages, neovascularisation, cotton-wool spots and microaneurysms) on a crop of the retinal region. For each of the 27 test images the crop is located in the original image by a coordinate-descent search over the four edges of the box, started from our FOV box, under two resampling hypotheses (the box stretched to the square, or padded to a square and resized uniformly), and accepted only when the normalised cross-correlation between the resampled original and the release reaches 0.97 (all 27: padded, 0.9988--0.9996). The official masks and the FOV are resampled through the recovered geometry with nearest-neighbour interpolation; cotton-wool spots are read as soft exudates and the other refined classes are ignored. Each IDRiD-trained checkpoint's original-resolution test maps are moved through the same geometry and scored by the protocol evaluator against both label sets, with the thresholds of its own validation run; no model is retrained or re-run. The per-class differences between the two label sets (pixels, components, intersection over union) are reported with the scores.

\subsection{Statistics}
Confidence intervals are 95\% percentile intervals from 2,000 bootstrap replicates over test images (seed 0); each replicate pools the histograms of the drawn images and recomputes AUPR, and a replicate without positive pixels is dropped; the interval of \mAUPR{} averages the four class values of the same replicates and drops a replicate in which any class has no positive pixel. Point estimates use the histograms at full 16-bit resolution and bootstrap replicates use per-image histograms coarsened to 4,096 bins. Two models are compared by a paired bootstrap on AUPR with shared image indices, and by a two-sided Wilcoxon signed-rank test on per-image Dice with Holm correction. Seed mean and s.d.\ are reported separately from bootstrap intervals. Rankings are compared across test sets with the Friedman test and the Nemenyi critical difference\cite{demsar2006statistical},
\begin{equation}
\mathrm{CD}=q_{\alpha}\sqrt{\frac{k(k+1)}{6N}} ,
\end{equation}
for $k$ models and $N=8$ test sets (two internal and six independent). Scores are the seed means of Table~\ref{tab:main}; in an independent block the score of a model is the mean over its IDRiD-trained and DDR-trained checkpoints; the DiaRetDB1 block uses the official image-level AUC on DiaRetDB1 (mean over the four classes). A version restricted to the fine-grained test sets is given in the repository's analysis notebook; the distribution of ranks under the image bootstrap\cite{maierhein2018rankings} is left for the final version. Agreement between published and retrained rankings is measured by Kendall's $\tau$.

\subsection{Time and resource consumption}
Training ran on NVIDIA A100 cards with 40\,GB in a node of eight, one task per card, with a queue that assigns tasks from longest to shortest expected duration. Evaluation of author weights, the first author-recipe rerun, data export and all metric computation ran on a workstation (RTX 4060 Ti, 16\,GB). Where a card cannot hold the effective batch of a recipe, gradients are accumulated and the deviation is recorded. Parameter count, floating-point operations, peak training memory, GPU hours and inference time per image are reported for every model (Extended Data Fig.~10).

\subsection{Data availability}
IDRiD, DDR and TJDR are public. FGADR and Retinal-Lesions are available from their authors on request. The FGADR licence forbids redistribution of the images and of any derived data; Retinal-Lesions is distributed under application terms, and we handle it in the same way. DiaRetDB1 is public. The MAPLES-DR labels are public (CC BY 4.0); the MESSIDOR images they annotate are available from ADCIS on request. We release all split lists, the protocol files and, for IDRiD, DDR, TJDR, MAPLES-DR and DiaRetDB1, the probability maps of every model (the maps hold no image data; DiaRetDB1 is distributed for research use). For FGADR and Retinal-Lesions we release split lists and code only.

\subsection{Code availability}
The benchmark code, including data export, training under the unified protocol, inference, the evaluation module and the scripts that produce every figure and table of this paper, is available at \url{https://github.com/skylynf/retiseg}. Official model repositories are referenced as pinned submodules and are not modified.

\FloatBarrier
\begin{table*}[t]
\centering\tablefont
\caption{\textbf{What each model keeps under the unified protocol.} Data, splits, the 1,440-pixel canvas, augmentation, the training amount in pixels, checkpoint selection and evaluation are shared (Methods). Loss, optimizer, learning rate, weight decay, schedule shape and effective batch are taken from each model's source recipe; the schedule is stretched over the shared number of steps. Models with a fixed input size train on patches of that size and predict with half-overlapping windows; all others train on $960\times960$ patches and predict on the whole canvas. nnU-Net runs through its own entry point and keeps its own plan.}
\label{tab:recipes}
\begin{adjustbox}{max width=\textwidth}
\begin{threeparttable}
\begin{tabular}{@{}l l l l l c c l c c l r@{}}
\toprule
\colh{Model} & \colh{Initialization} & \colh{Head} & \colh{Loss} & \colh{Optim.} & \colh{LR} & \colh{WD} & \colh{Schedule} & \colh{Batch} & \colh{Patch} & \colh{Inference} & \colh{Steps} \\
\midrule
\famchip{general}\; U-Net & none & 4 sigmoid & BCE + Dice$^a$ & Adam & 1e-4 & 0 & constant & 4 & 960 & whole canvas & 21,000 \\
\famchip{general}\; DeepLabv3 & ImageNet, ResNet-101 & 4 sigmoid & BCE & SGD & 7e-3 & 1e-4 & poly 0.9 & 16 & 960 & whole canvas & 6,000 \\
\famchip{general}\; HRNet-W48 & ImageNet & 4 sigmoid & BCE & SGD & 1e-2 & 5e-4 & poly & 12 & 960 & whole canvas & 7,000 \\
\famchip{general}\; Swin-Unet & ImageNet, Swin-T & 5 softmax & 0.4 CE + 0.6 Dice & SGD & 1e-2 & 1e-4 & poly 0.9 & 24 & 224 & sliding window & 63,000 \\
\famchip{general}\; FCT & none & 4 sigmoid & BCE & Adam & 1e-3 & 0 & constant & 4 & 384 & sliding window & 129,000 \\
\famchip{auto}\; nnU-Net & none & 5 softmax & CE + batch Dice & SGD & 1e-2 & 3e-5 & poly 0.9 & 2 & 1024$\times$1536$^c$ & sliding window & 250,000$^c$ \\
\famchip{auto}\; nnU-Net ResEnc-M & none & 5 softmax & CE + batch Dice & SGD & 1e-2 & 3e-5 & poly 0.9 & 13$^c$ & 384$\times$640$^c$ & sliding window & 250,000$^c$ \\
\famchip{fundus}\; M2MRF & ImageNet, HRNet-W48 & 4 sigmoid & binary Dice & SGD & 1e-2 & 5e-4 & poly 0.9 & 4 & 960 & whole canvas & 21,000 \\
\famchip{fundus}\; HACDR-Net & none & 5 softmax & weighted CE$^b$ & AdamW & 6e-5 & 1e-2 & poly 1.0, warm-up & 1 & 960 & whole canvas & 83,000 \\
\famchip{fundus}\; HRDecoder & ImageNet, HRNet-W48 & 4 sigmoid & binary Dice + 0.1 Dice (HR crops) & SGD & 1e-2 & 5e-4 & poly 0.9 & 4 & 960 & sliding window & 21,000 \\
\famchip{fundus}\; PMCNet & ImageNet, EfficientNet-B0 & 5 softmax & CE & Adam & 1e-4 & 0 & poly 0.9 & 1 & 960 & whole canvas & 83,000 \\
\famchip{fundus}\; WSRFNet & ImageNet, EfficientNet-B1 & 4 sigmoid & binary Dice, per loop$^d$ & SGD & 1e-2 & 5e-4 & poly 0.9 & 4 & 960 & whole canvas & 21,000 \\
\famchip{fundus}\; CLC-Net$^{\rm C}$ & none & 5 softmax & Dice + BCE (cls.\ heads)$^e$ & Adam & 3e-4 & 0 & step (1/3, 5/6) & 16 & 512 & sliding window & 73,000 \\
\famchip{fundus}\; TC-Net$^{\rm C}$ & none & 5 softmax & CE & Adam & 5e-4 & 1e-5 & constant & 4 & 512 & sliding window & 73,000 \\
\famchip{foundation}\; SAMed & SAM ViT-B, LoRA r=4 & 5 softmax & 0.2 CE + 0.8 Dice & AdamW & 5e-3 & 0.1 & warm-up, poly & 24 & 512 & sliding window & 13,000 \\
\famchip{foundation}\; SAM-Adapter & SAM ViT-H, adapters & 4 sigmoid & BCE + IoU & AdamW & 2e-4 & 1e-2 & stepped cosine & 2 & 1,024 & sliding window & 37,000 \\
\famchip{foundation}\; SAM2-UNet & SAM 2 Hiera-L, adapters & 4 sigmoid & BCE + weighted IoU$^f$ & AdamW & 1e-3 & 5e-4 & stepped cosine & 12 & 960 & whole canvas & 7,000 \\
\famchip{foundation}\; MedSAM (full FT) & MedSAM ViT-B, all weights & 4 sigmoid & Dice + BCE & AdamW & 1e-4 & 1e-2 & constant & 2 & 1,024 & sliding window & 37,000 \\
\famchip{foundation}\; RETFound + UPerNet & RETFound ViT-L (MAE) & 4 sigmoid & BCE + 0.4 aux. & AdamW & 4.7e-4 & 0.05 & warm-up, cosine & 24 & 224 & sliding window & 63,000 \\
\bottomrule
\end{tabular}
\begin{tablenotes}[flushleft]\scriptsize
\item Batch is the effective batch (gradient accumulation where one card cannot hold it). Steps is for IDRiD at 1,000 equivalent epochs; one equivalent epoch draws as many patch pixels as the training canvases contain, so steps = $\lceil\lceil$canvas pixels\,/\,patch pixels$\rceil$\,/\,batch$\rceil$ per epoch. DDR uses 400 equivalent epochs (for example 83,600 steps at batch 4 and 960-pixel patches; Supplementary Table~3 lists every run). $^{\rm C}$, re-implemented from the paper. Every run's recipe record, with its deviations from the source recipe, is in the repository.
\item $^a$Per-image, per-class sigmoid Dice with smoothing 1 in the registered recipe; the main-table runs pool the Dice term over the batch (Methods). The other Dice losses use $10^{-5}$. $^b$IDRiD: class-weighted cross-entropy; DDR: the same plus the authors' NALoss with weight 0.2. $^c$nnU-Net plans patch and batch per dataset (IDRiD plans shown; DDR: default $1024\times1536$ at batch 2, ResEnc-M $512\times512$ at batch 12) and trains 1,000 epochs of 250 steps on both sources; its learning rate, weight decay and polynomial schedule are its defaults. $^d$The paper's loss and optimizer, written into the wrapper because the repository has no training loop. $^e$Local and contextual Dice plus binary cross-entropy of the two classification heads, as described in the paper. $^f$Registered per-image term; the main-table run pools it over the batch (Methods).
\end{tablenotes}
\end{threeparttable}
\end{adjustbox}
\end{table*}

